%document class
\documentclass[10pt,oneside]{book}
%%%% Page Info + Commands %%%%%{

%packages
\usepackage{pgfplots}
\pgfplotsset{compat=1.18}
\usepackage{amsmath}
\usepackage{amsfonts}
\usepackage{amsthm,letltxmacro}
\usepackage{amssymb}
\usepackage{enumerate}
\usepackage{enumitem}
\usepackage[dvipsnames]{xcolor}
\usepackage{changepage}

%This will shrink the page
\newcommand{\smallpage}{  \setlength \oddsidemargin{.5in}
            \setlength \textwidth{5in}
            \setlength \topmargin{0in}
            \setlength \textheight{9in}}


% Commands for sets of numbers
\newcommand{\Z}{\mathbb{Z}}\newcommand{\R}{\mathbb{R}}\newcommand{\N}{\mathbb{N}}

% gordie spacing and boxing commands
\newcommand{\gspc}{\vspace{0.13cm}} % 0.5 space
\newcommand{\gline}{\gspc\\} % 1.5 space
\newcommand{\gbline}{\gspc\gspc\\} % double space
\newcommand{\gbox}[1]{\fbox{\parbox{\linewidth}{#1}}} % multiline box command


\newcommand{\gimage}[3]{
\begin{figure}[h]   
    \centering          
    \includegraphics[width=#2\textwidth]{#1}  
    \caption{#3}
    \label{fig:#1}
\end{figure}\\
}

\newcommand{\centerit}[1]{{\begin{center} #1 \end{center}}}
\newcommand{\ul}[1]{\underline{#1}}

%%%%%%%% command for graphics %%%%%%%%%%%%%
\usepackage{fancyhdr}
%}

%%%% Page Setup %%%%%%%
\smallpage\sloppy
\pagestyle{fancy}
\lhead{\large{\textbf{Day 3 - Berliner lectures us hooray}}}
%\rfoot{\thepage/\pageref{LastPage} }
\setlength{\headheight}{10pt}


% Proof writing
\LetLtxMacro\oldproof\proof
\let\endoldproof\endproof
\renewenvironment{proof}[1][\proofname]
  {\vspace{0.2cm}\begin{adjustwidth}{2em}{2em}
   \oldproof[#1]}
  {\endoldproof
   \end{adjustwidth}}



\begin{document}

\newcommand{\cg}[1]{\color{OliveGreen}#1\color{white}}
\newcommand{\cb}[1]{\color{BlueViolet}#1\color{white}}


%
% Introductions
%
\subsection*{Last Time in Class}
Last class, we primarily covered:
\begin{itemize}[itemsep=0pt, parsep=0pt]
	\item Divisibility
	\item GCD (greatest common denominator)
	\item Euclidean Algorithm
\end{itemize}
\subsubsection*{Euclidean Algorithm Review}
Below is an example with the Euclidean algorithm:
\begin{adjustwidth}{2em}{2em}\gspc
	\textit{EX:} gcd$(2346, 1002)$ 
	\begin{align*}
		2346 &= 2\cdot 1002 + 342&=\text{gcd}(2346, 1002)\\
		1002 &= 2\cdot 342 + 318&=\text{gcd}(1002, 342)\\
		342 &= 2\cdot 342 + 318&=\text{gcd}(342, 318)\\
		318 &= 13\cdot 24 + 6&=\text{gcd}(318, 24)\\
		24 &= 4\cdot 6 + 0&=\text{gcd}(24, 6)
	\end{align*}
	Finally, we see that our gcd is 6. 
\end{adjustwidth}\gspc
Berliner tells us that we likely will \textit{not have to do this many steps on our midterm exam}, so don't stress about it too much.
%
% GCD Review
%
\subsubsection*{GCD Theorem Review}
Note that the below theorem is quite important to remember:
\begin{adjustwidth}{2em}{2em}\gspc
	\textbf{Theorem!} There exist integers $x,y$ such that for integers $a,b$,
	\[\text{gcd}(a,b)=xa+yb.\]
	This theorem in plain english states that there is a linear combination of $a$ and $b$ that equals the greatest common divisor. \gspc
\end{adjustwidth}
This theorem has some pretty interesting implications. Take our previous example with gcd$(2346, 1002)$. We can work backworks from our gcd of 6 to find this linear combination:
\begin{align*}
	\cg {\fbox{6}} &= \fbox{318} - 32\cdot\fbox{24}\\
	\cg{\fbox{6}} &= \fbox{318} - 13\left[\fbox{342} - 1\cdot \fbox{318}\right] = -13\cdot\fbox{342} + 14\cdot \fbox{318}\\
	\cg{\fbox{6}} &= -32\cdot\fbox{342}+14\cdot\left[\fbox{1002}-2\cdot\fbox{342}\right]=14\cdot\fbox{1002}-41\cdot\fbox{342}\\
	\cg{\fbox{6}}&=14\cdot\fbox{1002}-41\cdot\left[\fbox{2346}-2\cdot\fbox{1002}\right]=\cb{\boxed{-41\cdot\fbox{2346}+96\cdot\fbox{1002}}}.
\end{align*}
Thus we've found the linear combination that gets the gcd (-41 and 96). What a painful process but it works. 
\pagebreak
\section*{Solving Diophantine Equations}
Recall from the last lecture that Diophantine Equations are equations with integer coefficients and integer solutions. For example, with the equation $x + y =1$, the solutions are all integers $x$ such that $x = 1 -y$. \textit{Now let's do an example}. 
\begin{adjustwidth}{2em}{2em}\gspc
	\textit{Solve.} $2346x + 1002y = 6$.\gline
	One possible solution: $x = -41, \;y=96$ (which we discussed above)
\end{adjustwidth}
However, solving for these integer solutions is not always so easy.
\subsection*{Fermat's Last Theorem}
Take for example, Fermat's Last Theorem:
\[x^n+y^n = z^n,\;\;x,y,z,n\in\Z\]
Notice that this equationresembles the form of a Pythagorean Triple $a^2 + b^2 = c^2$. However, the interesting thing is, for non-trivial solutions (e.g. $n = 0, 1, 2$, or $x = -y$ for odd $n$), that \textit{no solutions exist.}\gline
For now, we will primarily cover the equation $ax+ by=c$
%
%
%
\subsection*{Zesty Warm-Up}
Let's do some fun warm up problems:
\begin{enumerate}
	\item [(a)] What is the gcd of two consecutive even integers?\gline
	Consider that we have two integers, $n$ and $n + 2$. We can represent $n$ as $n = 2k$ for $k\in \Z$. Thus, we write:
	\begin{align*}
		d&\mid 2k\\
		d&\mid 2k+2
	\end{align*}
	Thus we have an $x,y\in \Z$ such that:
	\begin{align*}
		xd &= 2k\\
		yd &= 2k+ 2
	\end{align*}
	Pluggin in gives us that:
	\begin{align*}
		yd &= xd + 2\\
		(y-x)d &= 2
	\end{align*}
	Thus, for integer $d$, $d$ must equal -2, -1, 1, 2.\gline
	We can also try this with the euclidean algorithm variation;
	\begin{align*}
		2k + 2 &= 1(2k) + 2\\
		2k &= k(2) + 0\\
	\end{align*}
	Thus, we find that the gcd is the remainder one above 0, which is 2!
	% Second question
	\item [(b)]What is the gcd of two consecutive odd integers?
	Consider two consecutive odd integers, $n$ and $n +2$. We know that $n = 2k + 1$, for $k\in \Z$, so we can write:
	\begin{align*}
		d&\mid 2k + 1\\
		d&\mid 2k + 3
	\end{align*}
	We know that there exists integers $x,y$ such that:
	\begin{align*}
		xd &= 2k + 1\\
		yd &= 2k + 3
	\end{align*}
	Substituion yields:
	\begin{align*}
		xd &= yd + 2\\
		(x-y)d &= 2
	\end{align*}
	However, we know that $d$ cannot be 2 because $n$ and $n+2$ are odd, so $d = -1, 1$ and the gcd is $1$.\gline
	We also use the Euclidean algorithm to solve this problem:
	\begin{align*}
		2k +3 &= 1(2k+1) + 2\\
		2k+1 &= k(2) + \fbox{1}\\
		2 &= 2(1) + 0
	\end{align*}
	Thus, we see that the gcd is 1. 
\end{enumerate}
Here's the moral of this story: for gcd problems, the Euclidean algorithm is more consistent and more reliable than using the defintiion of a divisor. 
\subsection*{Diophantine Equation Practice}
Consider the Diophantine Equation $152x + 108y = 84$.
\begin{enumerate}
	\item [(a)] Show that $x = -3$, $y = 5$ is a solution:
	\begin{align*}
		152(-3) + 108(5) = 84
	\end{align*}
	\item [(b)] Can you find a solution for which $x > 0$?\gline
	Yes, $x = 108-3$ and $y = 152 -5$
	\item [(c, d)] Are there solutions where $x$ and $y$ are of the same sign?\gline
	No, there are no solutions that satisfy this condition 
\end{enumerate}

\subsection*{Another Diophantine Equation}
\begin{enumerate}
	\item Consider the Diophantine equation $6x + 9y = c$.
	\begin{enumerate}
		\item Are there any solutions for $x,y$ when $c = 6$?\gline
		Yes, when $x = 1$ and $y = 0$. 
		\item Are there any solutions for $x,y$ when $c = 3$?\gline
		Yes, when $x = -1$ and $y = 1$.\gline
		\item Are there any solutions for $x,y$ when $c = 33$?\gline
		Yes, when $x = 8$ and $y = -5$
		\item Key Takeaway: \textbf{Diophantine equations of the form $ax + by = c$ only have solutions if $c$ is a multiple of $\text{gcd}(a, b)$}
		\item What is the \textit{smallest} positive value of $c$ for which there is a solution for $x, y$?\gline
		It's 3. 
	\end{enumerate}
	\item Consider the Diophantine Equation $ax + by =1 $
	\begin{enumerate}
		\item Does this equation have any solutions when $(a,b) = (6,9)?$\gline
		No, because 6 and 9 are not relatively prime.
		\item Are there any solutions when $(a,b) = (2,4)$?\gline
		No, because 2 and 4 are not relatively prime.
		\item Are there any solutions when $(a,b) = (2,5)$?\gline
		Yes, because 2 and 5 are relatively prime.
		\item Are there any solutions when $(a,b) = (2k, 2k+1)$?
		Yes, because two consecutive integers are always relatively prime
		\item \textbf{Group Fill-in Theorem}:\\
		The diophantine equation $ax + by = 1$ has a soution for $(x,y)$ if and only if \ul{$a$ and $b$ are relatively prime.}
	\end{enumerate}
\end{enumerate}

% Section
\subsection*{Takeaways from the practice problems}
So what are we supposed to take away from these exercises? It's this one primary detail:
\begin{adjustwidth}{2em}{2em}\gspc\gspc
	\gbox {
	\textbf{Theorem! }
	The Diophantine eq. $ax+by=c$ only has a solution when $\text{gcd}(a,b) \mid c$}
\end{adjustwidth}\gspc\gspc
For example, with the equation $6x + 9y = c$, we know that we will have solutions as long as $c = \{\ldots, -3, 0, 3, 6, \ldots\}$. However, as soon as $3 \nmid c$, like with $1, 4, 10, -2$, \ul{no solution exists}. 
\gbline
However, what \ul{are} the solutions to $ax + by = c$? To get \ul{one} $(x_0, y_0)$, use Euclid Alg to get gcd($a,b$), then multiply by $\frac{c}{\text{gcd}(a,b)}.$ \gline
What about the others though?\gline
Pretend that $(x_1, y_1)$ is a solution, also. Then, we have that:
\begin{align*}
	\begin{cases}
		ax_0 + by_0 &= c\\
		ax_1 + by_1 &= c
	\end{cases}
\end{align*}
So setting both equations equal gives us that $ax_0 + by_0 = ax_1 + by_1$. This is quite helpful, as we can rewrite this as:
\begin{align*}
	a(x_1 - x_0) = b(y_0 - y_1)
\end{align*}
However, here's the little bitty crucial component. Remember that because $c \mid\text{gcd}(a,b)=d$, we have that $d\mid a$, and $d\mid b$. Thus, $\exists K, L\in \Z$ such that $Kd = a$ and $Ld = b$. Soooooo... look at this:
\begin{align*}
	Kd(x_1 - x_0) &= Ld(y_0 -y_1)\\
	K(x_1-x_0) &= L(y_0-y_1)
\end{align*}
Interestingly, a consequence of $d$ being the gcd is that $\text{gcd}(K,L) = 1$ (they are relatively prime). \\
\footnotesize *For the reason for why this is, I'll probably write a brief appendix to these notes but I don't have time to detail it right now. \normalsize\\
So, via the properties of divisibility, we know that:
\begin{align*}
	L&\mid (x_1 - x_0)\\
	K&\mid (y_0 - y_1) 
\end{align*}
Thus, there exists a $t\in \Z$ such that $Kt = (y_0 - y_1)$.
\pagebreak
So, we then write:
\begin{align*}
	K(x_1-x_9)&= L(Kt)\\
	x_1-x_0 &= Lt\\
	x_1 -x_0 &= Lt\\
	x_1 - x_0 &= \left(\frac{b}d\right) t\\
	x_1 &= x_0 + \left(\frac{b}d\right) t
\end{align*}
Furthermore, we can write this in terms of $y$ too:
\begin{align*}
	y_1 = y_0 - \left(\frac{a}d\right)t
\end{align*}
If you want both equations in a big box so they're easy to see, look no futher:
\begin{gather*}
	\cb{\boxed{x_1 = x_0 + \left(\frac{b}d\right) t}}\\
	\cb{\boxed{y_1 = y_0 - \left(\frac{a}d\right)t}}
\end{gather*}
\section*{Algorithmically Solving Diophantine Equations}
Thus, we estbalish a consistent algorithm for solving diophantine equations:\gspc\\
\gbox{
	\begin{enumerate}[topsep=0pt, itemsep=-1pt, parsep=-1pt]
		\item Use Euclidean Algorithm to solve your equation to get the gcd.
		\item Multiply by your $c$.
		\item Use the equations to solve your equations.
	\end{enumerate}
}\gspc\\\\
For example, let's consider the $6x + 9y = c$ equation from earlier:
\begin{enumerate}[itemsep=0pt, parsep=1pt]
	\item Use Euclidean algorithm\gline
	$6x + 9y = 3\ldots x = -1, y = 1$
	\item Multiply by 33:\gline
	$x_0 = -33, y_0=33$
	\item Solutions are:\gline
	$$x_1 = -33 + \frac{9}3 t = -33 +3t$$
	$$y_1 = 33 - \frac{6}3t = 33 - 2t$$ 
\end{enumerate}
Hooray! Wasn't that great? That's all the notes for now. 

\end{document}
